\documentclass{article}
\usepackage[UTF8]{ctex}
\usepackage{amsmath,amssymb,amsfonts}
\usepackage{enumitem}
\usepackage{geometry}
\geometry{a4paper, margin=1in}
% === 错题本排版优化（不参与题目标识） ===
\setlength{\parindent}{0pt}
\setlength{\parskip}{0.35em plus 0.10em minus 0.08em}
\setlength{\emergencystretch}{3em}
\setlist{itemsep=0.25em, topsep=0.40em, parsep=0pt}
\allowdisplaybreaks[2]
% === 错题本排版优化结束 ===
\begin{document}
\section*{中国海洋大学2026年数学分析真题}
\section*{126．中国海洋大学2026年数学分析试卷（回忆版）}

\subsection{题目1}
1．计算题．

\begin{itemize}
  \item[（1）] 求极限 \(\lim _{n \rightarrow \infty}\left(\ln n-\sum_{k=1}^{n} \frac{\ln k}{n}\right)\) ．
  \item[（2）] 求极限 \(\lim _{x \rightarrow 0} \frac{(1+x)^{x}-\cos \frac{\pi}{2}}{x\left(\sin x-\sin \frac{\pi}{2}\right)}\) ．
  \item[（3）] 求积分 \(\int_{0}^{1} \frac{x^{3}-x}{\ln x} \mathrm{~d} x\) ，需说明计算的合理性．
\end{itemize}

\subsection{题目2}
2．判断题．正确的给出证明，错误的给出反例．

\begin{itemize}
  \item[（1）] 任何数列都有单调的子列。
  \item[（2）] 若级数 \(\sum_{n=1}^{\infty}\left(a_{3 n-2}+a_{3 n-1}+a_{3 n}\right)\) 收敛，且 \(\lim _{n \rightarrow \infty} a_{n}=0\) ，则 \(\sum_{n=1}^{\infty} a_{n}\) 收敛。
  \item[（3）] 设 \(f(x)\) 在 \((-\infty,+\infty)\) 上二阶可导，则对任意的 \(\xi \in(-\infty,+\infty)\) ，一定存在 \(x_{1}, x_{2} \in(-\infty,+\infty)\) ，满足 \(f^{\prime}(\xi)=\frac{f\left(x_{2}\right)-f\left(x_{1}\right)}{x_{2}-x_{1}}\) ．
  \item[（4）] 若 \(f(x)\) 在 \([a,+\infty)\) 上单调非负且无界，则 \(\int_{a}^{+\infty} f(x) \mathrm{d} x\) 必发散．
\end{itemize}

\subsection{题目3}
3．解答如下问题：

\begin{itemize}
  \item[（1）] 若 \(\lim _{n \rightarrow \infty} \frac{a_{1}+a_{2}+\cdots+a_{n}}{n}=a\) ，则 \(\lim _{n \rightarrow \infty} a_{n}=a\) 。
  \item[（2）] 若 \(f(x)\) 在 \([a,+\infty)\) 上可导，且 \(\lim _{x \rightarrow a^{+}} f(x)=\lim _{x \rightarrow+\infty} f(x)=A\)（ \(A\) 为有限数），则存在 \(\xi \in(a,+\infty)\) ，使得 \(f^{\prime}(\xi)=0\) ．

\subsection{题目4}
\end{itemize}

\subsection{题目4}
原题边界“4．”：证明函数项级数 \(\sum_{n=1}^{\infty} a_{n} \frac{1}{n!} \int_{0}^{x} e^{-t} t^{n} \mathrm{~d} t\) 一致收敛的充要条件为级数 \(\sum_{n=1}^{\infty} a_{n}\) 收敛。

\subsection{题目5}
5．设 \(z=z(x, y)\) 二阶连续可微，求 \(a\) 的值，使得通过 \(u=x+3 y, v=x+a y\) 将 \(6 z_{x x}+z_{x y}-z_{y y}=0\)化为 \(z_{u v}=0\) 。

\subsection{题目6}
6．解答如下问题：

\begin{itemize}
  \item[（1）] 证明级数 \(\sum_{n=1}^{\infty} \frac{(-1)^{n}}{x^{n}+\frac{2^{n}}{n(n+1)}}\) 在 \((-\infty,+\infty)\) 上一致收敛。
  \item[（2）] 求 \(\lim _{x \rightarrow 0} \sum_{n=1}^{\infty} \frac{(-1)^{n}}{x^{n}+\frac{2^{n}}{n(n+1)}}\) ．
\end{itemize}

\subsection{题目7}
7．计算 \(\iint_{D} \sqrt{\left|y-x^{2}\right|} \mathrm{d} x \mathrm{~d} y\) ，其中 \(D=\left\{(x, y) \mid x^{2} \leq y \leq 2\right\}\) 。
\end{document}
